\documentclass[a4paper]{article}

% Basic packages
\usepackage[fontset=none]{ctex}
\usepackage{amsmath, amssymb}
\usepackage{graphicx}
\usepackage[margin=2.4cm]{geometry}
\usepackage[most]{tcolorbox}
\usepackage{etoolbox}
\usepackage{listings}
\usepackage{hyperref}
\usepackage{xurl}
\usepackage{booktabs}
\usepackage{longtable}
\usepackage{tabularx}
\usepackage{array}
\usepackage{subcaption}
\usepackage{float}
\usepackage{enumitem}
\usepackage{tikz}
\IfFileExists{pgfplots.sty}{
    \usepackage{pgfplots}
    \pgfplotsset{compat=1.18}
}{}

% Prefer LXGW fonts when available; fall back to TeX Live defaults otherwise.
\IfFontExistsTF{LXGW WenKai}{
    \setmainfont{LXGW WenKai}
    \setsansfont{LXGW WenKai}
    \setCJKmainfont{LXGW WenKai}
    \setCJKsansfont{LXGW WenKai}
}{
    \IfFontExistsTF{LXGW WenKai GB Screen}{
        \setmainfont{LXGW WenKai GB Screen}
        \setsansfont{LXGW WenKai GB Screen}
        \setCJKmainfont{LXGW WenKai GB Screen}
        \setCJKsansfont{LXGW WenKai GB Screen}
    }{
        \IfFontExistsTF{霞鹜文楷}{
            \setmainfont{霞鹜文楷}
            \setsansfont{霞鹜文楷}
            \setCJKmainfont{霞鹜文楷}
            \setCJKsansfont{霞鹜文楷}
        }{
            \setmainfont{Liberation Serif}
            \setsansfont{Liberation Sans}
            \setCJKmainfont{Noto Sans CJK JP}
            \setCJKsansfont{Noto Sans CJK JP}
        }
    }
}
\IfFontExistsTF{LXGW WenKai Mono}{
    \setmonofont{LXGW WenKai Mono}
    \setCJKmonofont{LXGW WenKai Mono}
}{
    \setmonofont{Noto Sans Mono}
}

% Hyperlinks
\hypersetup{
    colorlinks=true,
    linkcolor=blue!50!black,
    urlcolor=blue!60!black,
    citecolor=blue!50!black
}

% Highlight boxes
\newtcolorbox{findingbox}[1]{
    enhanced,
    colback=green!5!white,
    colframe=green!45!black,
    colbacktitle=green!45!black,
    coltitle=white,
    fonttitle=\bfseries,
    title=#1,
    sharp corners
}

\newtcolorbox{evidencebox}[1]{
    enhanced,
    colback=blue!4!white,
    colframe=blue!65!black,
    colbacktitle=blue!65!black,
    coltitle=white,
    fonttitle=\bfseries,
    title=#1,
    sharp corners
}

\newtcolorbox{riskbox}[1]{
    enhanced,
    colback=red!4!white,
    colframe=red!70!black,
    colbacktitle=red!70!black,
    coltitle=white,
    fonttitle=\bfseries,
    title=#1,
    sharp corners
}

\newtcolorbox{methodbox}[1]{
    enhanced,
    colback=yellow!9!white,
    colframe=yellow!60!black,
    colbacktitle=yellow!60!black,
    coltitle=black,
    fonttitle=\bfseries,
    title=#1,
    sharp corners
}

% Listings
\lstset{
    language=Python,
    basicstyle=\ttfamily\small,
    keywordstyle=\color{blue},
    stringstyle=\color{red!60!black},
    commentstyle=\color{green!50!black},
    breaklines=true,
    frame=single,
    numbers=left,
    numberstyle=\tiny\color{gray},
    captionpos=b,
    extendedchars=false
}

% Table helpers
\newcolumntype{Y}{>{\raggedright\arraybackslash}X}
\setlist[itemize]{leftmargin=1.8em}
\renewcommand{\arraystretch}{1.18}
\setlength{\tabcolsep}{6pt}

% Metadata
\newcommand{\reporttitle}{[在此填写研究主题]}
\newcommand{\reportsubtitle}{[在此填写副标题或研究切口]}
\newcommand{\reportauthors}{Codex}
\newcommand{\reportdate}{\today}
\newcommand{\researchquestion}{[在此填写核心研究问题]}
\newcommand{\reportscope}{[在此填写研究范围、时间边界或限定条件]}
\newcommand{\reportaudience}{[在此填写目标读者]}
\newcommand{\sourcewindow}{[在此填写资料截止日期或证据窗口]}
\newcommand{\reportcoverpath}{}

\begin{document}

\begin{titlepage}
\centering
{\Large 深度研究报告\par}
\vspace{1.0cm}
{\huge\bfseries \reporttitle\par}
\vspace{0.5cm}
{\Large \reportsubtitle\par}
\vspace{0.9cm}
{\large \reportauthors\par}
\vspace{0.2cm}
{\large \reportdate\par}
\vspace{1.0cm}

\ifdefempty{\reportcoverpath}{
{\small 若有合适的封面图、关键图表、产品截图或示意图，请在生成时填入本地路径。\par}
}{
\includegraphics[width=0.84\textwidth,height=0.42\textheight,keepaspectratio]{\reportcoverpath}\par
}

\vfill
\begin{tcolorbox}[width=0.94\textwidth, colback=black!2!white, colframe=black!60, sharp corners]
\textbf{核心问题}：\researchquestion\par
\textbf{研究范围}：\reportscope\par
\textbf{目标读者}：\reportaudience\par
\textbf{证据窗口}：\sourcewindow\par
\end{tcolorbox}
\end{titlepage}

\tableofcontents
\newpage

%% --- 正文内容开始 --- %%

% Recommended structure:
%
% \section{执行摘要}
% - Give the main conclusion first.
% - Separate strong findings, likely findings, and open uncertainties.
%
% \section{研究范围与方法}
% - Explain the question, scope, cutoff date, and evidence inventory.
% - Use methodbox or evidencebox when reproducibility matters.
%
% \section{主体分析}
% - Organize by thesis, comparison axis, or timeline rather than source order.
% - Use figures, tables, formulas, and citations whenever they improve understanding.
%
% \section{反例、风险与争议点}
% - Explicitly surface contradictory evidence or unresolved uncertainty.
%
% \section{结论与后续问题}
% - Distinguish confirmed facts, careful inference, and next validation steps.
%
% Suggested figure provenance pattern:
% \begin{figure}[H]
% \centering
% \includegraphics[width=\textwidth]{figures/example.png}
% \caption{关键图示 \protect\footnotemark}
% \end{figure}
% \footnotetext{来源：某网页 / 某 PDF 第 12 页 / 某视频 00:12:31--00:12:46。若为裁剪或重绘，请明确说明。}
%
% Suggested formula pattern:
% $$L = \sum_{i=1}^{n} w_i x_i$$
% \begin{itemize}
% \item $L$：总得分
% \item $w_i$：第 $i$ 项权重
% \item $x_i$：第 $i$ 项观测值
% \end{itemize}
%
% Suggested source appendix:
% \appendix
% \section{来源清单}
% \begin{longtable}{p{2cm}p{2cm}p{4cm}p{5cm}}
% \toprule
% Source ID & 模态 & 标题 / 标签 & 链接或路径 \\
% \midrule
% S1 & Web & 官方文档页面 & \url{https://example.com} \\
% S2 & Video & 演讲视频 & \url{https://example.com/video} \\
% \bottomrule
% \end{longtable}
%
% Post-build visual check for every self-drawn figure and every table:
% - no overflow, clipping, overlap, or caption collisions
% - no tiny labels or unreadable annotations
% - column widths and row heights remain comfortable to read
% - if a figure or table looks cramped, revise the layout and rebuild

%% --- 正文内容结束 --- %%

\end{document}
